\textbf{Phần 1. Trắc nghiệm lựa chọn}
\setcounterpageref{ex}{0}

\begin{ex} %[2D2V11]
Tính giá trị của phép tính: $13 - 5$.
\choice
{$7$}
{$9$}
{\True $8$}
{$6$}
\loigiai{
$13 - 5 = 8$.
\textbf{Đáp án của câu hỏi là C}
}
\end{ex}

\begin{ex} %[2D2V11]
Một đàn chim có $15$ con đậu trên dây điện. Lần thứ nhất có $6$ con bay đi, lần thứ hai có $3$ con bay đi. Hỏi còn lại bao nhiêu con chim trên dây điện?
\choice
{$7$ con chim}
{\True $6$ con chim}
{$5$ con chim}
{$8$ con chim}
\loigiai{
Số chim còn lại: $15 - 6 - 3 = 9 - 3 = 6$ (con chim).
\textbf{Đáp án của câu hỏi là B}
}
\end{ex}

\begin{ex} %[2D2V11]
Tìm $x$ biết: $12 - x = 4$.
\choice
{$7$}
{$6$}
{$9$}
{\True $8$}
\loigiai{
$x = 12 - 4 = 8$.
\textbf{Đáp án của câu hỏi là D}
}
\end{ex}

\begin{ex} %[2D2V11]
Năm nay Nam $8$ tuổi, anh trai Nam $14$ tuổi. Hỏi anh trai hơn Nam bao nhiêu tuổi?
\choice
{\True $6$ tuổi}
{$7$ tuổi}
{$5$ tuổi}
{$8$ tuổi}
\loigiai{
Anh hơn Nam: $14 - 8 = 6$ (tuổi).
\textbf{Đáp án của câu hỏi là A}
}
\end{ex}

\begin{ex} %[2D2V11]
Hiệu của $16 - 7$ bằng kết quả của phép tính nào dưới đây?
\choice
{$15 - 7$}
{$13 - 5$}
{\True $14 - 5$}
{$12 - 4$}
\loigiai{
$16 - 7 = 9$. Ta có $14 - 5 = 9$. Do đó hai phép tính có cùng kết quả.
\textbf{Đáp án của câu hỏi là C}
}
\end{ex}

\begin{ex} %[2D2V11]
Điền số thích hợp: $17 - 9 = 15 - \dots$
\choice
{$6$}
{\True $7$}
{$8$}
{$9$}
\loigiai{
Vế trái: $17 - 9 = 8$. Ta có: $15 - 8 = 7$. Vậy điền số 7.
\textbf{Đáp án của câu hỏi là B}
}
\end{ex}

\begin{ex} %[2D2V11]
Một sợi dây dài $15\text{ dm}$, người ta cắt bớt đi $8\text{ dm}$. Đoạn dây còn lại dài bao nhiêu?
\choice
{$6\text{ dm}$}
{$8\text{ dm}$}
{$9\text{ dm}$}
{\True $7\text{ dm}$}
\loigiai{
Đoạn dây còn lại dài: $15 - 8 = 7\text{ (dm)}$.
\textbf{Đáp án của câu hỏi là D}
}
\end{ex}

\begin{ex} %[2D2V11]
Tìm số tự nhiên biết rằng lấy $11$ trừ đi số đó thì được số bé nhất có một chữ số khác 0.
\choice
{\True $10$}
{$9$}
{$8$}
{$1$}
\loigiai{
Số bé nhất có một chữ số khác 0 là 1. Số đó là: $11 - 1 = 10$.
\textbf{Đáp án của câu hỏi là A}
}
\end{ex}

\begin{ex} %[2D2V11]
Một cửa hàng có $14$ chiếc quạt máy. Buổi sáng bán được $5$ chiếc, buổi chiều bán được $4$ chiếc. Cửa hàng còn lại bao nhiêu chiếc quạt?
\choice
{$6$ chiếc}
{$4$ chiếc}
{\True $5$ chiếc}
{$7$ chiếc}
\loigiai{
Tổng số quạt đã bán: $5 + 4 = 9$ (chiếc).\nSố quạt còn lại: $14 - 9 = 5$ (chiếc).
\textbf{Đáp án của câu hỏi là C}
}
\end{ex}

\begin{ex} %[2D2V11]
Trong một phép trừ, số bị trừ là $13$, hiệu là $6$. Số trừ là
\choice
{$6$}
{\True $7$}
{$8$}
{$5$}
\loigiai{
Số trừ = Số bị trừ - Hiệu = $13 - 6 = 7$.
\textbf{Đáp án của câu hỏi là B}
}
\end{ex}

\begin{ex} %[2D2V11]
Điền dấu thích hợp: $18 - 9 \dots 11 - 3$.
\choice
{$<$}
{$=$}
{Không so sánh được}
{\True $>$}
\loigiai{
$18 - 9 = 9$; $11 - 3 = 8$. Vì $9 > 8$ nên điền dấu $>$. 
\textbf{Đáp án của câu hỏi là D}
}
\end{ex}

\begin{ex} %[2D2V11]
Bác thợ có một thanh sắt dài $16\text{ dm}$. Bác cắt làm hai đoạn, đoạn thứ nhất dài $7\text{ dm}$. Hỏi đoạn thứ hai dài hơn đoạn thứ nhất bao nhiêu đề-xi-mét?
\choice
{\True $2\text{ dm}$}
{$9\text{ dm}$}
{$1\text{ dm}$}
{$3\text{ dm}$}
\loigiai{
Đoạn thứ hai dài: $16 - 7 = 9\text{ (dm)}$.\nĐoạn thứ hai hơn đoạn thứ nhất: $9 - 7 = 2\text{ (dm)}$.
\textbf{Đáp án của câu hỏi là A}
}
\end{ex}

\textbf{Phần 2. Trắc nghiệm đúng sai}
\setcounterpageref{ex}{0}

\begin{ex} %[2D2V11]
Xét tính đúng sai của các phép tính trong bảng trừ có nhớ:
\choiceTF
{\True Phép trừ $12 - 7$ có kết quả bằng $5$.}
{Phép trừ $13 - 6$ có kết quả lớn hơn $14 - 6$.}
{\True Khi cùng bớt ở số bị trừ và số trừ đi 2 đơn vị thì hiệu không thay đổi.}
{\True Trong phạm vi 20, phép trừ $11 - 9$ có hiệu nhỏ nhất trong các phép trừ có nhớ với số bị trừ 11.}
\loigiai{
\begin{itemize}
    \item Ý a) Đúng: $12 - 7 = 5$.
    \item Ý b) Sai: $13 - 6 = 7 < 14 - 6 = 8$.
    \item Ý c) Đúng: Tính chất của phép trừ.
    \item Ý d) Đúng: Vì 9 là số trừ lớn nhất có 1 chữ số nên $11 - 9 = 2$ là nhỏ nhất.
\end{itemize}
\begin{itemize}
    \item a) \textbf{Đúng}.
    \item b) \textbf{Sai}.
    \item c) \textbf{Đúng}.
    \item d) \textbf{Đúng}.
\end{itemize}
}
\end{ex}

\begin{ex} %[2D2V11]
Có 14 quyển vở chia cho hai bạn Nam và Việt. Nam nhận 8 quyển vở:
\choiceTF
{\True Việt nhận được 6 quyển vở.}
{\True Nam nhận nhiều hơn Việt 2 quyển vở.}
{Nếu Nam cho Việt 2 quyển thì hai bạn có số vở bằng nhau.}
{Tổng số vở của hai bạn là một số lẻ.}
\loigiai{
\begin{itemize}
    \item Ý a) Đúng: $14 - 8 = 6$ quyển.
    \item Ý b) Đúng: $8 - 6 = 2$ quyển.
    \item Ý c) Sai: Nam chỉ hơn Việt 2 quyển nên Nam chỉ cần cho Việt $2 : 2 = 1$ quyển là bằng nhau.
    \item Ý d) Sai: $14$ là số chẵn.
\end{itemize}
\begin{itemize}
    \item a) \textbf{Đúng}.
    \item b) \textbf{Đúng}.
    \item c) \textbf{Sai}.
    \item d) \textbf{Sai}.
\end{itemize}
}
\end{ex}

\textbf{Phần 3. Trả lời ngắn}
\setcounterpageref{ex}{0}

\begin{ex} %[2D2V11]
Tính giá trị của biểu thức: $15 - 9 + 8$.
\shortans{14}
\loigiai{
\textbf{Đáp số:} $14$\\
$15 - 9 = 6$; $6 + 8 = 14$.
}
\end{ex}

\begin{ex} %[2D2V11]
Tìm số tự nhiên x biết: $14 - x = 9$.
\shortans{5}
\loigiai{
\textbf{Đáp số:} $5$\\
$x = 14 - 9 = 5$.
}
\end{ex}

\begin{ex} %[2D2V11]
Trong hộp có 13 viên bi gồm bi xanh và bi đỏ. Biết số bi đỏ là 6 viên. Hỏi số bi xanh nhiều hơn số bi đỏ bao nhiêu viên?
\shortans{1}
\loigiai{
\textbf{Đáp số:} $1$\\
Số bi xanh: $13 - 6 = 7$ viên. Hơn bi đỏ: $7 - 6 = 1$ viên.
}
\end{ex}

\begin{ex} %[2D2V11]
Một sợi dây dài 17 cm, bạn Mai cắt đi một đoạn thì còn lại đoạn dây dài 9 cm. Hỏi đoạn dây bị cắt đi dài bao nhiêu xăng-ti-mét?
\shortans{8}
\loigiai{
\textbf{Đáp số:} $8$\\
Đoạn dây bị cắt: $17 - 9 = 8\text{ cm}$.
}
\end{ex}

\begin{ex} %[2D2V11]
Tìm số bị trừ biết số trừ là 8 và hiệu là 7.
\shortans{15}
\loigiai{
\textbf{Đáp số:} $15$\\
Số bị trừ = $8 + 7 = 15$.
}
\end{ex}

\begin{ex} %[2D2V11]
Lớp 2A có 16 bạn tham gia múa hát, trong đó có 9 bạn nữ. Hỏi số bạn nam ít hơn số bạn nữ bao nhiêu bạn?
\shortans{2}
\loigiai{
\textbf{Đáp số:} $2$\\
Số bạn nam: $16 - 9 = 7$ bạn. Nam ít hơn nữ: $9 - 7 = 2$ bạn.
}
\end{ex}
